% BHU PYQ -- Political Science: Principles of Political Science-II (2024-25 NEP)
\documentclass[11pt,a4paper]{article}
\usepackage[a4paper,top=2.5cm,bottom=2.5cm,left=2.8cm,right=2.8cm,headheight=14pt]{geometry}
\usepackage{amsmath,amssymb}
\usepackage{fontspec}
\setmainfont{Kohinoor Devanagari}
\usepackage{microtype,fancyhdr,enumitem,hyperref,lastpage,parskip}

\hypersetup{colorlinks=true,urlcolor=black,linkcolor=black,pdftitle={POLMJ21: Principles of Political Science-II -- BHU NEP 2024-25}}

\pagestyle{fancy}\fancyhf{}
\renewcommand{\headrulewidth}{0.4pt}\renewcommand{\footrulewidth}{0.4pt}
\fancyhead[L]{\small \textit{Banaras Hindu University}}
\fancyhead[R]{\small \textit{POLMJ21: Political Science-II (2024-25)}}
\fancyfoot[C]{\small Page \thepage\ of \pageref{LastPage}}

\newlist{parts}{enumerate}{2}
\setlist[parts,1]{label=(\alph*),leftmargin=*,itemsep=4pt,topsep=3pt}
\setlist[parts,2]{label=(\roman*),leftmargin=*,itemsep=2pt,topsep=2pt}

\newcommand{\pts}[1]{\hfill \textit{[#1]}}

\begin{document}

\begin{center}
  {\large\bfseries BANARAS HINDU UNIVERSITY}\\[2pt]
  {\normalsize B. A. (Hons.) Semester II Examination, 2024-25 (NEP)}\\[6pt]
  \rule{0.8\linewidth}{0.5pt}\\[6pt]
  {\large\bfseries POLITICAL SCIENCE}\\[4pt]
  {\large\bfseries Paper Code: POLMJ-21 : Principles of Political Science-II}\\[4pt]
  {\normalsize Time: 2:30 Hours \hfill Full Marks: 70}\\[2pt]
  {\small REF NO. CE/Conf./D-II/01}
\end{center}
\rule{\linewidth}{0.4pt}
\smallskip

\noindent \textit{Instructions: This question paper comprises three Sections A, B and C. All questions are compulsory. Write your Roll No.\ at the top immediately on receipt of this question paper.}

\smallskip
\noindent \textit{इस प्रश्न-पत्र में तीन खण्ड अ, ब तथा स हैं। सभी प्रश्न अनिवार्य हैं।}

\bigskip

\section*{SECTION -- A / खण्ड -- अ}
\noindent \textit{Note: Answer each question within 50 words. Each question carries 2 marks.}\pts{5 $\times$ 2 = 10}

\noindent \textit{प्रत्येक प्रश्न का उत्तर 50 शब्दों के भीतर दीजिए। प्रत्येक प्रश्न 2 अंकों का है।}

\begin{enumerate}[label=\textbf{\arabic*.},leftmargin=*,itemsep=8pt]
  \item 
  \begin{parts}
    \item Negative Liberty / नकारात्मक स्वतन्त्रता
    \item Procedural Justice / प्रक्रियात्मक न्याय
    \item Direct Democracy / प्रत्यक्ष लोकतन्त्र
    \item Natural Rights / प्राकृतिक अधिकार
    \item Self-Determination / आत्म-निर्णय
  \end{parts}
\end{enumerate}

\bigskip

\section*{SECTION -- B / खण्ड -- ब}
\noindent \textit{Note: Answer the following questions in about 250 words. Each question carries 10 marks.}\pts{3 $\times$ 10 = 30}

\noindent \textit{निम्नलिखित प्रश्नों का उत्तर लगभग 250 शब्दों में दीजिए। प्रत्येक प्रश्न 10 अंकों का है।}

\begin{enumerate}[label=\textbf{\arabic*.},leftmargin=*,start=2,itemsep=12pt]

  \item Explain Amartya Sen's notion on Equality.\pts{10}
  
  समानता के संदर्भ में अमर्त्य सेन के विचारों की व्याख्या कीजिए।

  \begin{center}\textbf{OR / अथवा}\end{center}

  What is Democracy? What are its major variants?\pts{10}
  
  लोकतन्त्र क्या है? इसके मुख्य प्रकार क्या हैं?

  \item Define Nationalism and explain its core values.\pts{10}
  
  राष्ट्रवाद को परिभाषित कीजिए तथा इसके मुख्य मूल्यों की व्याख्या कीजिए।

  \begin{center}\textbf{OR / अथवा}\end{center}

  Differentiate between Nationalism and Patriotism.\pts{10}
  
  राष्ट्रवाद तथा देशभक्ति के मध्य अन्तर स्पष्ट कीजिए।

  \item Describe the development of the concept of Citizenship during different historical periods.\pts{10}
  
  विभिन्न कालखण्डों के दौरान नागरिकता की अवधारणा के विकास का वर्णन कीजिए।

  \begin{center}\textbf{OR / अथवा}\end{center}

  Elaborate J. S. Mill's notion on Liberty.\pts{10}
  
  स्वतन्त्रता के संदर्भ में जे० एस० मिल के विचारों की व्याख्या कीजिए।

\end{enumerate}

\bigskip

\section*{SECTION -- C / खण्ड -- स}
\noindent \textit{Note: Answer all the questions. Each question carries 15 marks.}\pts{2 $\times$ 15 = 30}

\noindent \textit{सभी प्रश्नों के उत्तर दीजिए। प्रत्येक प्रश्न 15 अंकों का है।}

\begin{enumerate}[label=\textbf{\arabic*.},leftmargin=*,start=5,itemsep=14pt]

  \item Describe different perspectives of different thinkers on Equality.\pts{15}
  
  समानता पर विभिन्न विचारकों के अलग-अलग दृष्टिकोणों का वर्णन कीजिए।

  \begin{center}\textbf{OR / अथवा}\end{center}

  Explain John Rawls' theory of Justice.\pts{15}
  
  जॉन रॉल्स के न्याय के सिद्धान्त की व्याख्या कीजिए।

  \item Describe different theories of Rights.\pts{15}
  
  अधिकारों के विभिन्न सिद्धान्तों का वर्णन कीजिए।

  \begin{center}\textbf{OR / अथवा}\end{center}

  Explain the difference between Negative and Positive Liberty with suitable examples.\pts{15}
  
  उपयुक्त उदाहरणों की सहायता से नकारात्मक तथा सकारात्मक स्वतन्त्रता के मध्य अन्तर स्पष्ट कीजिए।

\end{enumerate}

\end{document}
